\newif\ifglobal

%%%%%%%%%%%%%%%%%%%%%%%
%%% Compile to view %%%
%%%%%%%%%%%%%%%%%%%%%%%

%%%%%%%%%%%%%%%%%%%%%%%%%%%%%%%%%%%%%%%%%%%%%%%%%%%
%%% Readme:                                     %%%
%%% https://www.overleaf.com/read/bwdjvfjksvjy/ %%%
%%%%%%%%%%%%%%%%%%%%%%%%%%%%%%%%%%%%%%%%%%%%%%%%%%%

% >>> M?: marks a module setting number or important notice

% >>> M4: Set {./relative-path} to external files for this module <<<
\def \modulefiles {./23-SYSREG}
% To add an external file, use:
% (Replace '---' with actual filename)
% '\input{\modulefiles/tables/---.tex}' for tables
% '\includegraphics{\modulefiles/figures/---}' for figures
% '\subfile{\modulefiles/---}' for register files

% >>> M1: This module is in global project? <<<
% 'YES' - uncomment; 'NO' - comment out
\globaltrue

\ifglobal
    \documentclass[main\_\_EN.tex]{subfiles}

\else
    \documentclass[openany, 10pt]{book}

    % >>> M2: In ./00-shared/preamble-trm-module, also find
    %         settings for visibility of todo notes, watermark <<<
    \usepackage{preamble-trm-module}

    % >>> M3: Temporary labels in English,
    %         you can add your own labels there <<<
    \myexternaldocument{./00-shared/temp-labels-en}

\fi %\ifglobal


\setcounter{secnumdepth}{4}

% >>> M5: If this module needs any updates in
% './00-shared/preamble-trm-module.sty'
% (include additional package, etc.), contact your module PM
% or create the file './00-shared/changelog.tex' make the
% necessary updates yourself and document them in this file <<<


\begin{document}


% Sets language into which pre-defined bits of text are translated
\selectlanguage{english}

% Do not delete this block
\ifglobal
% Add the following front matter if
% this module is outside of global project
\else
    \InputIfFileExists{readme}{}{}
    {\let\clearpage\relax \listoftodos}
    \tableofcontents
    %\listoftables
    %\listoffigures
\fi


%%%%%%%%%%%%%%%%%%%%%%%%%
%%% TRM Module Begins %%%
%%%%%%%%%%%%%%%%%%%%%%%%%

\hypertarget{sysreg}{}
\chapter{System Registers (SYSTEM)}\label{mod:sysreg}

\section{Overview}

The \chipname{} integrates a large number of peripherals, and enables the control of individual peripherals to achieve optimal characteristics in performance-vs-power-consumption scenarios. Specifically, \chipname{} has a various of system configuration registers that can be used for the chip's clock management (clock gating), power management, and the configuration of peripherals and core-system modules. This chapter lists all these system registers and their functions.

\section{Features}

\chipname{} system registers can be used to control the following peripheral blocks and core modules:

\begin{itemize}
    \item System and memory
    \item Clock
    \item Software Interrupt
    \item Low-power management
    \item Peripheral clock gating and reset
    \item CPU Control
\end{itemize}


\section{Function Description}

\subsection{System and Memory Registers}

\subsubsection{Internal Memory}

The following registers can be used to control \chipname{}'s internal memory:

\begin{itemize}
\item In register \hyperref[regdescSYSCONCLKGATEFORCEONREG]{SYSCON\_CLKGATE\_FORCE\_ON\_REG}:
    \begin{itemize}
        \item Setting different bits of the \hyperref[fielddescSYSCONROMCLKGATEFORCEON]{SYSCON\_ROM\_CLKGATE\_FORCE\_ON} field forces on the clock gates of different blocks of Internal ROM 0 and Internal ROM 1.
        \item Setting different bits of the \hyperref[fielddescSYSCONSRAMCLKGATEFORCEON]{SYSCON\_SRAM\_CLKGATE\_FORCE\_ON} field forces on the clock gates of different blocks of Internal SRAM.
        \item This means when the respective bits of this register are set to 1, the clock gate of the corresponding ROM or SRAM blocks will always be on. Otherwise, the clock gate will turn on automatically when the corresponding ROM or SRAM blocks are accessed and turn off automatically when the corresponding ROM or SRAM blocks are not accessed. Therefore, it's recommended to configure these bits to 0 to lower power consumption.
    \end{itemize}
    \item In register \hyperref[regdescSYSCONMEMPOWERDOWNREG]{SYSCON\_MEM\_POWER\_DOWN\_REG}:
    \begin{itemize}
        \item Setting different bits of the \hyperref[fielddescSYSCONROMPOWERDOWN]{SYSCON\_ROM\_POWER\_DOWN} field sends different blocks of Internal ROM 0 and Internal ROM 1 into retention state.
        \item Setting different bits of the \hyperref[fielddescSYSCONSRAMPOWERDOWN]{SYSCON\_SRAM\_POWER\_DOWN} field sends different blocks of Internal SRAM into retention state.
        \item The ``Retention" state is a low-power state of a memory block. In this state, the memory block still holds all the data stored but cannot be accessed, thus reducing the power consumption. Therefore, you can send a certain block of memory into the retention state to reduce power consumption if you know you are not going to use such memory block for some time.
    \end{itemize}
    \item In register \hyperref[regdescSYSCONMEMPOWERUPREG]{SYSCON\_MEM\_POWER\_UP\_REG}:
    \begin{itemize}
        \item By default, all memory enters low-power state when the chip enters the Light-sleep mode.
        \item Setting different bits of the \hyperref[fielddescSYSCONROMPOWERUP]{SYSCON\_ROM\_POWER\_UP} field forces different blocks of Internal ROM 0 and Internal ROM 1 to work as normal (do not enter the retention state) when the chip enters Light-sleep.
        \item Setting different bits of the \hyperref[fielddescSYSCONSRAMPOWERUP]{SYSCON\_SRAM\_POWER\_UP} field forces different blocks of Internal SRAM to work as normal (do not enter the retention state) when the chip enters Light-sleep.
    \end{itemize}
\end{itemize}

For detailed information about the controlling bits of different blocks, please see Table \ref{tab:sysreg-mem-controlling-bit} below.

\begin{longtable}[c]{| l | l | l |l | l | l |}
\caption{Internal Memory Controlling Bit}
\label{tab:sysreg-mem-controlling-bit}
\\\hline
\rowcolor{lightgray}
Internal Memory & Lowest Address1 & Highest Address1 & Lowest Address2 & Highest Address2 & Controlling Bit \\ \hline
        Internal ROM 0  & 0x{}4000\_0000    & 0x{}4003\_FFFF  &- & -  & Bit0 \\\hline
        Internal ROM 1  & 0x{}4004\_0000    & 0x{}4004\_FFFF  &-  & - & Bit1\\\hline
        Internal ROM 2 & 0x{}4005\_0000    & 0x{}4005\_FFFF & 0x{}3FF0\_0000 & 0x{}3FF0\_FFFF& Bit2 \\\hline

        SRAM Block0  & 0x{}4037\_0000    & 0x{}4037\_3FFF     & -    & -     & Bit0 \\\hline
        SRAM Block1  & 0x{}4037\_4000    & 0x{}4037\_7FFF     & -    & -     & Bit1 \\\hline
        SRAM Block2  & 0x{}4037\_8000    & 0x{}4037\_FFFF     & 0x{}3FC8\_8000    & 0x{}3FC8\_FFFF     & Bit2 \\\hline
        SRAM Block3  & 0x{}4038\_0000    & 0x{}4038\_FFFF     & 0x{}3FC9\_0000    & 0x{}3FC9\_FFFF     & Bit3 \\\hline
        SRAM Block4  & 0x{}4039\_8000    & 0x{}4039\_FFFF     & 0x{}3FCA\_0000    & 0x{}3FCA\_FFFF     & Bit4 \\\hline
        SRAM Block5  & 0x{}403A\_C000    & 0x{}403A\_FFFF     & 0x{}3FCB\_C000    & 0x{}3FCB\_FFFF     & Bit5 \\\hline

        SRAM Block6  & 0x{}403B\_0000    & 0x{}403B\_FFFF     & 0x{}3FCC\_0000    & 0x{}3FCC\_FFFF     & Bit6 \\\hline
        SRAM Block7  & 0x{}403C\_0000    & 0x{}403C\_FFFF     & 0x{}3FCD\_4000    & 0x{}3FCD\_FFFF     & Bit7 \\\hline
        SRAM Block8  & 0x{}403D\_0000    & 0x{}403D\_BFFF     & 0x{}3FCE\_8000    & 0x{}3FCE\_FFFF     & Bit8 \\\hline
        SRAM Block9  & -    & -     & 0x{}3FCF\_0000    & 0x{}3FCF\_7FFF     & Bit9 \\\hline

        SRAM Block10  & -    & -     & 0x{}3FCF\_8000    & 0x{}3FCF\_FFFF     & Bit10 \\\hline
\end{longtable}

For detailed information about the controlling bits of different blocks, please see Table \ref{tab:sysreg-mem-controlling-bit} below.

\subsubsection{External Memory}

\hyperref[regdesc:SYSTEMEXTERNALDEVICEENCRYPTDECRYPTCONTROLREG]{SYSTEM\_EXTERNAL\_DEVICE\_ENCRYPT\_DECRYPT\_CONTROL\_REG} configures encryption and decryption options of the external memory. For details, please refer to Chapter \ref{mod:extmemencr}
\textit{\nameref{mod:extmemencr}}.

\subsubsection{RSA Memory}

\hyperref[regdesc:SYSTEMRSAPDCTRLREG]{SYSTEM\_RSA\_PD\_CTRL\_REG} controls the SRAM memory in the RSA accelerator.

\begin{itemize}
    \item Setting the \hyperref[fielddesc:SYSTEMRSAMEMPD]{SYSTEM\_RSA\_MEM\_PD} bit to send the RSA memory into retention state. This bit has the lowest priority, meaning it can be masked by the \hyperref[fielddesc:SYSTEMRSAMEMFORCEPU]{SYSTEM\_RSA\_MEM\_FORCE\_PU} field. This bit is invalid when the \textit{\nameref{mod:digsig}} occupies the RSA.
    \item Setting the \hyperref[fielddesc:SYSTEMRSAMEMFORCEPU]{SYSTEM\_RSA\_MEM\_FORCE\_PU} bit to force the RSA memory to work as normal when the chip enters light sleep. This bit has the second highest priority, meaning it overrides the \hyperref[fielddesc:SYSTEMRSAMEMPD]{SYSTEM\_RSA\_MEM\_PD} field.
    \item Setting the \hyperref[fielddesc:SYSTEMRSAMEMFORCEPD]{SYSTEM\_RSA\_MEM\_FORCE\_PD} bit to send the RSA memory into retention state. This bit has the highest priority, meaning it sends the RSA memory into retention state regardless of the \hyperref[fielddesc:SYSTEMRSAMEMFORCEPU]{SYSTEM\_RSA\_MEM\_FORCE\_PU} field.
\end{itemize}

\subsection{Clock Registers}

The following registers are used to set clock sources and frequency. For more information, please refer to Chapter \ref{mod:resclk}
\textit{\nameref{mod:resclk}}.

\begin{itemize}
    \item \hyperref[regdesc:SYSTEMCPUPERCONFREG]{SYSTEM\_CPU\_PER\_CONF\_REG}
    \item \hyperref[regdesc:SYSTEMSYSCLKCONFREG]{SYSTEM\_SYSCLK\_CONF\_REG}
    \item \hyperref[regdesc:SYSTEMBTLPCKDIVFRACREG]{SYSTEM\_BT\_LPCK\_DIV\_FRAC\_REG}
\end{itemize}

\subsection{Interrupt Signal Registers}

The following registers are used for generating the interrupt signals, which then can be routed to the CPU peripheral interrupts via the interrupt matrix. To be more specific, writing 1 to any of the following registers generates an interrupt signal. Therefore, these registers can be used by software to control interrupts. For more information, please refer to Chapter \ref{mod:intmtrx}
\textit{\nameref{mod:intmtrx}}.

\begin{itemize}
    \item \hyperref[fielddesc:SYSTEMCPUINTRFROMCPU0]{SYSTEM\_CPU\_INTR\_FROM\_CPU\_0\_REG}
    \item \hyperref[fielddesc:SYSTEMCPUINTRFROMCPU1]{SYSTEM\_CPU\_INTR\_FROM\_CPU\_1\_REG}
    \item \hyperref[fielddesc:SYSTEMCPUINTRFROMCPU2]{SYSTEM\_CPU\_INTR\_FROM\_CPU\_2\_REG}
    \item \hyperref[fielddesc:SYSTEMCPUINTRFROMCPU3]{SYSTEM\_CPU\_INTR\_FROM\_CPU\_3\_REG}
\end{itemize}


\subsection{Low-power Management Registers}

The following registers are used for low-power management. For more information, please refer to Chapter \ref{mod:lowpowm} \textit{\nameref{mod:lowpowm}}.
\begin{itemize}
    \item \hyperref[regdesc:SYSTEMRTCFASTMEMCONFIGREG]{SYSTEM\_RTC\_FASTMEM\_CONFIG\_REG}: configures the RTC CRC check.
    \item \hyperref[regdesc:SYSTEMRTCFASTMEMCRCREG]{SYSTEM\_RTC\_FASTMEM\_CRC\_REG}: configures the CRC check value.
\end{itemize}

\subsection{Peripheral Clock Gating and Reset Registers}

The following registers are used for controlling the clock gating and reset of different peripherals. Details can be seen in Table \ref{tab:sysreg-peripheral-controlling-bit}.

\begin{itemize}
    \item \hyperref[regdesc:SYSTEMCACHECONTROLREG]{SYSTEM\_CACHE\_CONTROL\_REG}
    \item \hyperref[regdesc:SYSTEMEDMACTRLREG]{SYSTEM\_EDMA\_CTRL\_REG}
    \item \hyperref[regdesc:SYSTEMPERIPCLKEN0REG]{SYSTEM\_PERIP\_CLK\_EN0\_REG}
    \item \hyperref[regdesc:SYSTEMPERIPRSTEN0REG]{SYSTEM\_PERIP\_RST\_EN0\_REG}
    \item \hyperref[regdesc:SYSTEMPERIPCLKEN1REG]{SYSTEM\_PERIP\_CLK\_EN1\_REG}
    \item \hyperref[regdesc:SYSTEMPERIPRSTEN1REG]{SYSTEM\_PERIP\_RST\_EN1\_REG}

\end{itemize}

\chipname{} features low power consumption. This is why some peripheral clocks are gated (disabled) by default. Before using any of these peripherals, it is mandatory to enable the clock for the given peripheral and release the peripheral from reset state. For details, see the table below:

\begin{longtable}[c]{ |p{4cm}|p{6cm}|p{6cm}| }
\caption{Peripheral Clock Gating and Reset Bits}
\label{tab:sysreg-peripheral-controlling-bit}
\\\hline\rowcolor{lightgray}
Peripheral    & Clock Enabling Bit{\hyperref[other:sysreg-reset-notes]{$^1$}} & Reset Controlling Bit{\hyperref[other:sysreg-reset-notes]{$^2$}}{\hyperref[other:sysreg-reset-notes]{$^3$}}\\\hline
%%%%%%%%%%%%%%%%%%%%%%%%%%%%%%%%%%%%%%%%%%%%%%%%%%%%%%%%%%%%%%%%%%%%%%%%%%%%%%%%%%%%%%%%%%%%%%%%%%%%%%%%%%%%%%%%%%%%%%%%%%%

%%%%%%%%%%%%%%%%%%%%%%%%%%%%%%%%%%%%%%%%%%%%%%%%%%%%%%%%%%%%%%%%%%%%%%%%%%%%%%%%%%%%%%%%%%%%%%%%%%%%%%%%%%%%%%%%%%%%%%%%%%%
\textbf{EDMA Ctrl} & \multicolumn{2}{|c|}{\textbf{\hyperref[regdesc:SYSTEMEDMACTRLREG]{SYSTEM\_EDMA\_CTRL\_REG}}}\\\hline
EDMA & \hyperref[fielddesc:SYSTEMEDMACLKON]{SYSTEM\_EDMA\_CLK\_ON} & \hyperref[fielddesc:SYSTEMEDMARESET]{SYSTEM\_EDMA\_RESET}  \\\hline

\textbf{CACHE Ctrl} & \multicolumn{2}{|c|}{\textbf{\hyperref[regdesc:SYSTEMCACHECONTROLREG]{SYSTEM\_CACHE\_CONTROL\_REG}}}\\\hline
DCACHE & \hyperref[fielddesc:SYSTEMDCACHECLKON]{SYSTEM\_DCACHE\_CLK\_ON} & \hyperref[fielddesc:SYSTEMDCACHERESET]{SYSTEM\_DCACHE\_RESET} \\\hline
ICACHE & \hyperref[fielddesc:SYSTEMICACHECLKON]{SYSTEM\_ICACHE\_CLK\_ON} & \hyperref[fielddesc:SYSTEMICACHERESET]{SYSTEM\_ICACHE\_RESET} \\\hline

\textbf{Peripheral} & \textbf{\hyperref[regdesc:SYSTEMPERIPCLKEN0REG]{SYSTEM\_PERIP\_CLK\_EN0\_REG}} & \textbf{\hyperref[regdesc:SYSTEMPERIPRSTEN0REG]{SYSTEM\_PERIP\_RST\_EN0\_REG}} \\ \hline

Timer Group0    & \hyperref[fielddesc:SYSTEMTIMERGROUPCLKEN]{SYSTEM\_TIMERGROUP\_CLK\_EN} & \hyperref[fielddesc:SYSTEMTIMERGROUPRST]{SYSTEM\_TIMERGROUP\_RST} \\\hline
Timer Group1    & \hyperref[fielddesc:SYSTEMTIMERGROUP1CLKEN]{SYSTEM\_TIMERGROUP1\_CLK\_EN} & \hyperref[fielddesc:SYSTEMTIMERGROUP1RST]{SYSTEM\_TIMERGROUP1\_RST} \\\hline
System Timer    & \hyperref[fielddesc:SYSTEMSYSTIMERCLKEN]{SYSTEM\_SYSTIMER\_CLK\_EN} & \hyperref[fielddesc:SYSTEMSYSTIMERRST]{SYSTEM\_SYSTIMER\_RST} \\\hline
UART0    & \hyperref[fielddesc:SYSTEMUARTCLKEN]{SYSTEM\_UART\_CLK\_EN} & \hyperref[fielddesc:SYSTEMUARTRST]{SYSTEM\_UART\_RST} \\\hline
UART1    & \hyperref[fielddesc:SYSTEMUART1CLKEN]{SYSTEM\_UART1\_CLK\_EN} & \hyperref[fielddesc:SYSTEMUART1RST]{SYSTEM\_UART1\_RST} \\\hline
UART MEM    & \hyperref[fielddesc:SYSTEMUARTMEMCLKEN]{SYSTEM\_UART\_MEM\_CLK\_EN} {\hyperref[other:sysreg-reset-notes]{$^4$}} & \hyperref[fielddesc:SYSTEMUARTMEMRST]{SYSTEM\_UART\_MEM\_RST} \\\hline
SPI0, SPI1    & \hyperref[fielddesc:SYSTEMSPI01CLKEN]{SYSTEM\_SPI01\_CLK\_EN} & \hyperref[fielddesc:SYSTEMSPI01RST]{SYSTEM\_SPI01\_RST} \\\hline
SPI2    & \hyperref[fielddesc:SYSTEMSPI2CLKEN]{SYSTEM\_SPI2\_CLK\_EN} & \hyperref[fielddesc:SYSTEMSPI2RST]{SYSTEM\_SPI2\_RST} \\\hline
SPI3    & \hyperref[fielddesc:SYSTEMSPI3CLKEN]{SYSTEM\_SPI3\_CLK\_EN} & \hyperref[fielddesc:SYSTEMSPI3RST]{SYSTEM\_SPI3\_RST} \\\hline
I2C0    & \hyperref[fielddesc:SYSTEMI2CEXT0CLKEN]{SYSTEM\_I2C\_EXT0\_CLK\_EN} & \hyperref[fielddesc:SYSTEMI2CEXT0RST]{SYSTEM\_I2C\_EXT0\_RST} \\\hline
I2C1    & \hyperref[fielddesc:SYSTEMI2CEXT1CLKEN]{SYSTEM\_I2C\_EXT1\_CLK\_EN} & \hyperref[fielddesc:SYSTEMI2CEXT1RST]{SYSTEM\_I2C\_EXT1\_RST} \\\hline
I2S0    & \hyperref[fielddesc:SYSTEMI2S0CLKEN]{SYSTEM\_I2S0\_CLK\_EN} & \hyperref[fielddesc:SYSTEMI2S0RST]{SYSTEM\_I2S0\_RST} \\\hline
I2S1    & \hyperref[fielddesc:SYSTEMI2S1CLKEN]{SYSTEM\_I2S1\_CLK\_EN} & \hyperref[fielddesc:SYSTEMI2S1RST]{SYSTEM\_I2S1\_RST} \\\hline
TWAI Controller    & \hyperref[fielddesc:SYSTEMCANCLKEN]{SYSTEM\_CAN\_CLK\_EN} & \hyperref[fielddesc:SYSTEMCANRST]{SYSTEM\_CAN\_RST} \\\hline
UHCI0    & \hyperref[fielddesc:SYSTEMUHCI0CLKEN]{SYSTEM\_UHCI0\_CLK\_EN} & \hyperref[fielddesc:SYSTEMUHCI0RST]{SYSTEM\_UHCI0\_RST} \\\hline
USB    & \hyperref[fielddesc:SYSTEMUSBCLKEN]{SYSTEM\_USB\_CLK\_EN} & \hyperref[fielddesc:SYSTEMUSBRST]{SYSTEM\_USB\_RST} \\\hline
RMT    & \hyperref[fielddesc:SYSTEMRMTCLKEN]{SYSTEM\_RMT\_CLK\_EN} & \hyperref[fielddesc:SYSTEMRMTRST]{SYSTEM\_RMT\_RST} \\\hline
PCNT    & \hyperref[fielddesc:SYSTEMPCNTCLKEN]{SYSTEM\_PCNT\_CLK\_EN} & \hyperref[fielddesc:SYSTEMPCNTRST]{SYSTEM\_PCNT\_RST} \\\hline
PWM0    & \hyperref[fielddesc:SYSTEMPWM0CLKEN]{SYSTEM\_PWM0\_CLK\_EN} & \hyperref[fielddesc:SYSTEMPWM0RST]{SYSTEM\_PWM0\_RST} \\\hline
PWM1    & \hyperref[fielddesc:SYSTEMPWM1CLKEN]{SYSTEM\_PWM1\_CLK\_EN} & \hyperref[fielddesc:SYSTEMPWM1RST]{SYSTEM\_PWM1\_RST} \\\hline
LED\_PWM Controller    & \hyperref[fielddesc:SYSTEMLEDCCLKEN]{SYSTEM\_LEDC\_CLK\_EN} & \hyperref[fielddesc:SYSTEMLEDCRST]{SYSTEM\_LEDC\_RST} \\\hline
ADC Arbiter    & \hyperref[fielddesc:SYSTEMADC2ARBCLKEN]{SYSTEM\_ADC2\_ARB\_CLK\_EN} & \hyperref[fielddesc:SYSTEMADC2ARBRST]{SYSTEM\_ADC2\_ARB\_RST} \\\hline
ADC Controller    & \hyperref[fielddesc:SYSTEMAPBSARADCCLKEN]{SYSTEM\_APB\_SARADC\_CLK\_EN} & \hyperref[fielddesc:SYSTEMAPBSARADCRST]{SYSTEM\_APB\_SARADC\_RST} \\\hline
%%%%%%%%%%%%%%%%%%%%%%%%%%%%%%%%%%%%%%%%%%%%%%%%%%%%%%%%%%%%%%%%%%%%%%%%%%%%%%%%%%%%%%%%%%%%%%%%%%%%%%%%%%%%%%%%%%%%%%%%%%%
\textbf{Accelerators} & \textbf{\hyperref[regdesc:SYSTEMPERIPCLKEN1REG]{SYSTEM\_PERIP\_CLK\_EN1\_REG}} & \textbf{\hyperref[regdesc:SYSTEMPERIPRSTEN1REG]{SYSTEM\_PERIP\_RST\_EN1\_REG}} \\ \hline
USB\_DEVICE & \hyperref[fielddesc:SYSTEMUSBDEVICECLKEN]{SYSTEM\_USB\_DEVICE\_CLK\_EN} & \hyperref[fielddesc:SYSTEMUSBDEVICERST]{SYSTEM\_USB\_DEVICE\_RST} \\\hline
UART2 & \hyperref[fielddesc:SYSTEMUART2CLKEN]{SYSTEM\_UART2\_CLK\_EN} & \hyperref[fielddesc:SYSTEMUART2RST]{SYSTEM\_UART2\_RST} \\\hline
LCD\_CAM & \hyperref[fielddesc:SYSTEMLCDCAMCLKEN]{SYSTEM\_LCD\_CAM\_CLK\_EN} & \hyperref[fielddesc:SYSTEMLCDCAMRST]{SYSTEM\_LCD\_CAM\_RST} \\\hline
SDIO\_HOST & \hyperref[fielddesc:SYSTEMSDIOHOSTCLKEN]{SYSTEM\_SDIO\_HOST\_CLK\_EN} & \hyperref[fielddesc:SYSTEMSDIOHOSTRST]{SYSTEM\_SDIO\_HOST\_RST} \\\hline
DMA & \hyperref[fielddesc:SYSTEMDMACLKEN]{SYSTEM\_DMA\_CLK\_EN} & \hyperref[fielddesc:SYSTEMDMARST]{SYSTEM\_DMA\_RST} {\hyperref[other:sysreg-reset-notes]{$^5$}}  \\\hline
HMAC & \hyperref[fielddesc:SYSTEMCRYPTOHMACCLKEN]{SYSTEM\_CRYPTO\_HMAC\_CLK\_EN} & \hyperref[fielddesc:SYSTEMCRYPTOHMACRST]{SYSTEM\_CRYPTO\_HMAC\_RST} {\hyperref[other:sysreg-reset-notes]{$^6$}} \\\hline
Digital Signature & \hyperref[fielddesc:SYSTEMCRYPTODSCLKEN]{SYSTEM\_CRYPTO\_DS\_CLK\_EN} & \hyperref[fielddesc:SYSTEMCRYPTODSRST]{SYSTEM\_CRYPTO\_DS\_RST} {\hyperref[other:sysreg-reset-notes]{$^7$}} \\\hline
RSA Accelerator & \hyperref[fielddesc:SYSTEMCRYPTORSACLKEN]{SYSTEM\_CRYPTO\_RSA\_CLK\_EN} & \hyperref[fielddesc:SYSTEMCRYPTORSARST]{SYSTEM\_CRYPTO\_RSA\_RST} \\\hline
SHA Accelerator & \hyperref[fielddesc:SYSTEMCRYPTOSHACLKEN]{SYSTEM\_CRYPTO\_SHA\_CLK\_EN} & \hyperref[fielddesc:SYSTEMCRYPTOSHARST]{SYSTEM\_CRYPTO\_SHA\_RST} \\\hline
AES Accelerator & \hyperref[fielddesc:SYSTEMCRYPTOAESCLKEN]{SYSTEM\_CRYPTO\_AES\_CLK\_EN} & \hyperref[fielddesc:SYSTEMCRYPTOAESRST]{SYSTEM\_CRYPTO\_AES\_RST} \\\hline
peri backup & \hyperref[fielddesc:SYSTEMPERIBACKUPCLKEN]{SYSTEM\_PERI\_BACKUP\_CLK\_EN} & \hyperref[fielddesc:SYSTEMPERIBACKUPRST]{SYSTEM\_PERI\_BACKUP\_RST} \\\hline
\end{longtable}
\begin{tiplisting}
\vspace{-3em}\label{other:sysreg-reset-notes}
\begin{enumerate}
    \item Setting the clock enable register to 1 enables the clock, and to 0 disables the clock;
    \item Setting the reset enabling register to 1 resets a peripheral, and to 0 disables the reset.
    \item Reset registers cannot be cleared by hardware. Therefore, SW reset clear is required after setting the reset registers.
    \item UART memory is shared by all UART peripherals, meaning having any active UART peripherals will prevent the UART memory from entering the clock-gated state.
    \item When DMA is required for peripheral communications, for example, UCHI0, SPI, I2S, LCD\_CAM, AES, SHA and ADC, DMA clock should also be enabled.
    \item Resetting this bit also resets the SHA accelerator.
    \item Resetting this bit also resets the AES, SHA, and RSA accelerators.
\end{enumerate}
\end{tiplisting}

\subsection{CPU Control Registers}
These registers control CPU0 and CPU1 of \chipname. Note that, by default, only CPU0 is started when the SoC powers up. During this time, the clock of CPU1 is disabled. Therefore, users need to enable CPU1 clock manually to use both CPU0 and CPU1.
\begin{itemize}
    \item \hyperref[regdesc:SYSTEMCORE1CONTROL0REG]{SYSTEM\_CORE\_1\_CONTROL\_0\_REG}
    \begin{itemize}
        \item Setting the \hyperref[fielddesc:SYSTEMCONTROLCORE1RESETING]{SYSTEM\_CONTROL\_CORE\_1\_RESETING} bit resets CPU1.
    \item \hyperref[fielddesc:SYSTEMCONTROLCORE1CLKGATEEN]{SYSTEM\_CONTROL\_CORE\_1\_CLKGATE\_EN} controls the CPU1 clock.
    \item Setting the \hyperref[fielddesc:SYSTEMCONTROLCORE1RUNSTALL]{SYSTEM\_CONTROL\_CORE\_1\_RUNSTALL} bit stalls CPU1. When this bit is set, CPU1 will finish the on-going task and stalls.
    \end{itemize}
    \item Register \hyperref[regdesc:SYSTEMCORE1CONTROL1REG]{SYSTEM\_CORE\_1\_CONTROL\_1\_REG} is used to facilitate the communication between CPU0 and CPU1. To be more specific, one CPU can write this register, using the agreed-on formats, which will then be read by the other CPU, thus achieving communication between these two CPUs. Note that the value in this register will not affect any hardware configuration, which allows CPU communication solely controlled by software.
\end{itemize}


\clearpage
\section{Register Summary}
\hypertarget{sysreg-reg-summ}{}
\label{sec:sysreg-reg-summ}

In this section, the addresses of all the registers starting with SYSTEM are relative to the base address of system registers provided in Table \ref{tab:sysmem-base-address} in Chapter \ref{mod:sysmem} \textit{\nameref{mod:sysmem}}; and those starting with APB are relative to the base address of APB control also provided in Table \ref{tab:sysmem-base-address} in Chapter \ref{mod:sysmem} \textit{\nameref{mod:sysmem}}.

The abbreviations given in Column \textbf{Access} are explained in Section \hyperref[glossary-access-types]{\textit{Access Types for Registers}}.

\subfile{\modulefiles/23-SYSREG-reg-summ__EN}
\subfile{\modulefiles/23-SYSREG-apb-reg-summ__EN}


\clearpage
\section{Registers}

In this section, the addresses of all the registers starting with SYSTEM are relative to the base address of system registers provided in Table \ref{tab:sysmem-base-address} in Chapter \ref{mod:sysmem} \textit{\nameref{mod:sysmem}}; and those starting with APB are relative to the base address of APB control also provided in Table \ref{tab:sysmem-base-address} in Chapter \ref{mod:sysmem} \textit{\nameref{mod:sysmem}}.

\subfile{\modulefiles/23-SYSREG-reg__EN}
\subfile{\modulefiles/23-SYSREG-apb-reg__EN}

\end{document}
